\subsection{THE SCHEME OF SELECTION AND UNION}

The scheme of selection and union is the following :

S8. Let R be a relation, let x and y be distinct letters, and let X and Y be letters distinct from x and y which do not appear in R. Then the relation

\[
(1) \quad (\forall y) (\exists X) (\forall x) (R \Longrightarrow (x \in X)) \Longrightarrow (\forall Y) \operatorname{Coll} _ {X} ((\exists y) ((y \in Y) \text {   and   } R))
\]

is an axiom.

Let us first show that this rule is indeed a scheme. Let \(S\) denote the relation (1), and let us substitute a term \(T\) for a letter \(z\) in \(S\) ; by CS8 (Chapter I, §4, no. 1) we may suppose that \(x, y, X, Y\) are distinct from \(z\) and do not appear in \(T\) . Then \((T|z)S\) is identical with

\[
(\forall y) (\exists X) (\forall x) (R ^ {\prime} \Longrightarrow (x \in X)) \Longrightarrow (\forall Y) \operatorname{Coll} _ {X} ((\exists y) ((y \in Y) \text {   and   } R ^ {\prime})),
\]

where \(R'\) is \((T|z)R\) .

Intuitively, the relation \((\forall y)(\exists X)(\forall x)(R \Rightarrow (x \in X))\) means that for every object \(y\) there exists a set \(X\) (which may depend on \(y\) ) such that the objects \(x\) which are in the relation \(R\) with the given object \(y\) are elements of \(X\) (but not necessarily the whole of \(X\) ). The scheme of selection and union asserts that if this is the case and if \(Y\) is any set, then there exists a set whose elements are precisely the objects \(x\) which are in the relation \(R\) with at least one object \(y\) of the set \(Y\) .

C51. Let P be a relation, let A be a set, and let x be a letter which does not appear in A. Then the relation ``P and \(x \in A\) '' is collectivizing in x.

Let R denote the relation ``P and x = y'', where y is a letter distinct from x which appears neither in P nor in A. The relation

\[
(\forall x) (R \Longrightarrow (x \in \{y \}))
\]

is true by C27 (Chapter I, §4, no. 1). Let \(X\) be a letter distinct from \(x\) and \(y\) which does not appear in \(P\) . The preceding relation is identical with \((\{y\}|X)((\forall x)(R \Longrightarrow (x \in X)))\) (because \(x\) is distinct from \(y\) ), and therefore the relation \((\forall y)(\exists X)(\forall x)(R \Longrightarrow (x \in X))\) is true by virtue of S5 and C27 (Chapter I, §4, nos. 1 and 2). It follows from S8 and C30 (Chapter I, §4, no. 3) that the relation

\[
(A | Y) \operatorname{Coll} x (\exists y) (y \in Y \text {   and   } R)
\]

(where Y is a letter which does not appear in R) is true, and this relation is identical with \(\operatorname{Coll}_{\mathbf{x}}(\exists\mathbf{y})(\mathbf{y}\in\mathbf{A}\) and R) (because neither x nor y appears in A). Finally, the relation `` \(y\in A\) and R'' is equivalent to ``x=y and \(x\in A\) and P'' by C43 (Chapter I, §5, no. 1); since x appears neither in P nor in A, the relation

\[
(\exists y) (x = y \text {   and   } x \in A \text {   and   } P)
\]

is equivalent to ``((∃y)(x=y)) and x∈A and P'' by C33 (Chapter I, §4, no. 3) and therefore to ``P and x∈A'' because (∃y)(x=y) is true.

The set \(\mathcal{E}_{\pmb{x}}(\pmb{P}\) and \(\pmb{x} \in \pmb{A})\) is called the set of all \(\pmb{x} \in \pmb{A}\) such that \(\pmb{P}\) (* thus we may speak of the set of all real numbers such that \(\pmb{P}_{*}\) ).

C52. Let R be a relation, A a set, and x a letter which does not appear in A. If the relation \(R \Longrightarrow (x \in A)\) is a theorem, then R is collectivizing in x.

For R is then equivalent to ``R and \(x \in A\) ''.

Remark. Let R be a relation which is collectivizing in x, and let S be a relation such that \((\forall x)(S \Longrightarrow R)\) is a theorem. Then S is collectivizing in x; for R is equivalent to \(x \in \mathcal{E}_{x}(R)\) , so that

\[
S \Longrightarrow (x \in \mathcal {E} _ {x} (R))
\]

is a theorem, and the assertion follows from C52. Notice also that in this case we have \(\mathcal{E}_{\pmb{X}}(\pmb{S}) \subset \mathcal{E}_{\pmb{X}}(\pmb{R})\) by C50.

C53. Let T be a term, A a set, x and y distinct letters. Suppose that x does not appear in A and that y does not appear in A nor in T. Then the relation \((\exists x)(y = T \text{ and } x \in A)\) is collectivizing in y.

Let R be the relation y = T. The relation \((\forall y)(R \Longrightarrow (y \in \{T\}))\) is true, hence so is \((\forall x)(\exists X)(\forall y)(R \Longrightarrow (y \in X))\) , where X is a letter, distinct from y, which does not appear in R. By virtue of S8, the relation \((\exists x)(x \in A \text{ and } R)\) is collectivizing in y, and C53 is proved.

The relation \((\exists x)(y = T\) and \(x \in A)\) is often read as follows: ``y can be put in the form \(T\) for some \(x\) belonging to \(A\)''. The set

\[
\mathcal {E} _ {\boldsymbol {y}} ((\exists \boldsymbol {x}) (\boldsymbol {y} = \boldsymbol {T} \text {   and   } \boldsymbol {x} \in A))
\]

is generally called the set of objects of the form T for \(x \in A\) . The assembly so denoted contains neither x nor y, and does not depend on the choice of the letter y satisfying the conditions of C53.

